% TOP_PROBLEM: {"id": 30007067, "title": "Two-step linear images of special quasi-projective groups", "category_id": 5, "category": {"id": 5, "name": "algebraic_geometry", "display_name": "Algebraic Geometry", "description": "Geometric objects defined by polynomial equations.", "slug": "algebraic-geometry", "order_index": 5, "created_at": "2026-07-31T15:26:25.671Z"}, "status": "open"}
% FIRST_POSTED: 2026-10-06
% ATTEMPT_STATUS: solved
% OPENAI_CLAIM: {"schema_version":1,"source_repo":"https://github.com/openai/math","source_commit":"adc7f1241b42e322a6451854ab7e4b4c146bf78a","release_date":"2026-10-06","match":"full","resolution":"proved","independently_reviewed":false,"second_pass":"agreed","families":[{"family":"057","title":"Fundamental groups of special complex varieties and root orbifolds","subject":"Algebraic and complex geometry","manuscripts":[{"title":"Two-step monodromy of special quasi-projective varieties","pdf_path":"preprints/Two-step-monodromy-of-special-quasi-projective-varieties-September-24-2026/paper.pdf","readme_path":"preprints/Two-step-monodromy-of-special-quasi-projective-varieties-September-24-2026/README.md","date":"2026-09-24","theorem_ref":"preprints/Two-step-monodromy-of-special-quasi-projective-varieties-September-24-2026/build/sections/introduction.tex#thm:linear"}],"lean":null,"reasoning_trace":null}]}
\documentclass[11pt]{article}
\usepackage[margin=25mm]{geometry}
\usepackage{fontspec}
\setmainfont{Latin Modern Roman}
\usepackage{amsmath,amssymb,mathtools}
\usepackage{unicode-math}
\setmathfont{Latin Modern Math}
\usepackage{enumitem}
\usepackage{xurl,hyperref,bookmark}
\hypersetup{colorlinks=true,urlcolor=blue}
\setcounter{secnumdepth}{0}
\setlength{\parindent}{0pt}
\setlength{\parskip}{5pt}
\setlength{\emergencystretch}{3em}
\begin{document}
\section*{Two-step linear images of special quasi-projective groups}

\subsection{OpenAI's claim}
Every complex linear representation of a smooth special quasi-projective fundamental group is claimed to have virtually two-step-nilpotent image, while the paper also constructs nonspecial general quasi-Albanese fibers. The cited manuscript argues for that result; its relation to this problem is specified below.

\subsection{Outline of the argument}
Manuscript 1 (family 057, 2026-09-24): At each canonical nilpotent stage, the fiber Lie algebra equals the commutator algebra. Its abelianization has only weight minus two, making the next conjugation bracket incompatible with the allowed weights; universality then transfers the resulting class-two bound to arbitrary unipotent representations.

\subsection{Scope relative to this problem}
Manuscript 1 (family 057, 2026-09-24): The problem selects the mathematical statement summarized here, including its stated hypotheses and parameter range. The cited theorem gives a claimed proof of that assertion in full.

\subsection{Formal verification}
Family 057 has no Lean scope page in this snapshot; no formalization of these manuscript results is listed.

This site has not executed the Lean files or checked their dependencies. These descriptions report the scope of the linked declarations, rather than a new verification of the proofs.

\subsection{Status caveat}
OpenAI attributes these results to an internal model. Its release notice cautions that results lacking full formalization can contain mistakes. This site has not independently assessed the proofs; a solved label records the scope of the reported claim.

\subsection{Sources}
\begin{itemize}
\item OpenAI, \emph{Two-step monodromy of special quasi-projective varieties}, 2026-09-24; citation key \texttt{OAI:Two-step-monodromy-of-special-quasi-projective-varieties-September-24-2026}. \href{https://github.com/openai/math/blob/adc7f1241b42e322a6451854ab7e4b4c146bf78a/preprints/Two-step-monodromy-of-special-quasi-projective-varieties-September-24-2026/paper.pdf}{PDF (GitHub)}; \href{https://github.com/openai/math/raw/adc7f1241b42e322a6451854ab7e4b4c146bf78a/preprints/Two-step-monodromy-of-special-quasi-projective-varieties-September-24-2026/paper.pdf}{PDF (direct)}; \href{https://github.com/openai/math/blob/adc7f1241b42e322a6451854ab7e4b4c146bf78a/preprints/Two-step-monodromy-of-special-quasi-projective-varieties-September-24-2026/README.md}{Manuscript page and citation}; \href{https://github.com/openai/math/blob/main/preprints/Two-step-monodromy-of-special-quasi-projective-varieties-September-24-2026/paper.pdf}{Latest version}.
\item \href{https://github.com/openai/math/blob/adc7f1241b42e322a6451854ab7e4b4c146bf78a/CONTENTS.md}{OpenAI catalogue entry}.
\end{itemize}
\end{document}
